\chapter{Planar N=4 Supersymmetric Yang-Mills as a Spectral Problem}
\label{ch:planar-n4-spectral-problem-spin-chains}

\ConventionsLorentzian

The subject of this chapter is the spectrum of local gauge-invariant
operators in four-dimensional \(\mathcal N=4\) supersymmetric Yang--Mills
theory in the planar limit.  The same theory appears earlier in this volume
as a maximally supersymmetric conformal gauge theory.  Here it is treated as a
spectral problem: diagonalize the dilatation operator on the vector space of
renormalized local operators.

\section{Planar Limit and Operator Hilbert Space}

Let the gauge group be \(SU(N)\) or \(U(N)\), with gauge coupling defined by
the kinetic term
\[
  \frac{1}{4g_{\rm YM}^2}\int \dd^4x\,
  \operatorname{tr}(F_{\mu\nu}F^{\mu\nu})
\]
in the trace normalization used throughout the gauge-theory volumes.  The
't Hooft coupling is
\[
  \lambda=g_{\rm YM}^2N,\qquad
  g=\frac{\sqrt\lambda}{4\pi}.
\]
The planar limit means
\[
  N\to\infty,\qquad \lambda\ \text{fixed},
\]
with local single-trace operators normalized so that their planar two-point
functions remain finite.

Let \(\mathfrak V\) denote the one-letter vector space spanned by the
elementary adjoint fields of \(\mathcal N=4\) SYM and their covariant
derivatives, modulo the equations of motion and gauge-covariant descendant
relations appropriate to local-operator renormalization.  For fixed length
\(L\), a single-trace word is
\[
  \mathcal O_W(x)
  =
  \operatorname{tr}\!\bigl(W_1(x)W_2(x)\cdots W_L(x)\bigr),
  \qquad W_i\in\mathfrak V .
\]
Cyclicity of the trace identifies
\[
  W_1\otimes W_2\otimes\cdots\otimes W_L
  \sim
  W_2\otimes\cdots\otimes W_L\otimes W_1 .
\]
Thus the fixed-length single-trace space is the cyclic quotient
\[
  \mathcal H_L^{\rm cyc}
  =
  \mathfrak V^{\otimes L}/\langle\tau-1\rangle ,
\]
where \(\tau\) shifts all letters by one site.

The dilatation operator \(D\) is defined by the logarithmic scale dependence
of renormalized local operators.  In a basis \(\mathcal O_A\),
\[
  \mu\frac{\dd}{\dd\mu}\mathcal O_A
  =
  -\Gamma_A{}^B\mathcal O_B,\qquad
  D_A{}^B=(D_0)_A{}^B+\Gamma_A{}^B .
\]
At a conformal fixed point, diagonalizing \(D\) is equivalent to finding local
operators of definite scaling dimension.  Planarity implies that, at fixed
loop order, \(D\) acts locally along the cyclic word: at \(\ell\) loops the
range is at most \(\ell+1\) adjacent sites in the scalar subsectors described
below.

\section{The SU(2) Scalar Sector}

Choose two complex scalar fields
\[
  Z=\frac{1}{\sqrt2}(\phi^5+\ii\phi^6),
  \qquad
  X=\frac{1}{\sqrt2}(\phi^1+\ii\phi^2).
\]
The \(SU(2)\) scalar sector is spanned by
\[
  \operatorname{tr}(W_1\cdots W_L),\qquad W_i\in\{Z,X\}.
\]
Identify \(Z\) with a down spin and \(X\) with an up spin:
\[
  Z\leftrightarrow\ket{\downarrow},\qquad
  X\leftrightarrow\ket{\uparrow}.
\]
Let \(P_{j,j+1}\) be the operator that interchanges the two letters at sites
\(j\) and \(j+1\), with \(L+1\equiv1\).

\begin{theorem}[One-loop planar \(SU(2)\) Hamiltonian]
\label{thm:planar-n4-one-loop-su2}
In the \(SU(2)\) scalar sector, with the coupling convention above, the planar
one-loop dilatation operator is
\[
  D
  =
  L+\frac{\lambda}{8\pi^2}H_{XXX}
  +O(\lambda^2),
  \qquad
  H_{XXX}=\sum_{j=1}^{L}(1-P_{j,j+1})
\]
on the cyclic spin-chain quotient.
\end{theorem}

\begin{proof}
The only planar one-loop terms that move \(X\) and \(Z\) letters inside the
holomorphic \(SU(2)\) sector come from the scalar quartic interaction and the
wavefunction and non-exchange diagrams required by supersymmetry.  The
exchange part is local on adjacent sites and is proportional to
\(-P_{j,j+1}\).  The non-exchange part is proportional to the identity on the
same pair.  Hence the most general nearest-neighbor planar one-loop operator
compatible with the two-letter \(SU(2)\) symmetry and translation around the
trace is
\[
  c\sum_{j=1}^L(1-P_{j,j+1})+c'\sum_{j=1}^L(1+P_{j,j+1}).
\]
The half-BPS operators \(\operatorname{tr}(Z^L)\) and their fully symmetrized
scalar descendants have protected scaling dimension.  On a symmetric pair,
\(P_{j,j+1}=1\), so protection forces \(c'=0\).  The coefficient \(c\) is
fixed by the logarithmic part of the single adjacent exchange integral.  In
the present normalization the divergent scalar integral is
\[
  \frac{\lambda}{8\pi^2}\log(\mu |x|)+\text{finite},
\]
and its color contraction is planar and nearest-neighbor.  Substituting this
coefficient gives the displayed \(H_{XXX}\).  The cyclic quotient implements
the trace relation and imposes zero total spin-chain momentum on physical
single-trace states.
\end{proof}

The proof uses supersymmetric protection only to fix the identity term after
the exchange coefficient has been computed.  This is a useful place where the
field-theory origin of the spin chain matters: the spin-chain Hamiltonian is
not introduced independently of the renormalized local-operator problem.

\section{Coordinate Bethe Ansatz}

On an infinite chain, a one-magnon wavefunction is
\[
  \Psi_p(n)=\ee^{\ii pn}.
\]
The Hamiltonian \(H_{XXX}\) gives
\[
  (H_{XXX}\Psi_p)(n)
  =
  2\Psi_p(n)-\Psi_p(n+1)-\Psi_p(n-1),
\]
hence
\[
  E_1(p)=4\sin^2\frac p2 .
\]
The corresponding one-loop anomalous dimension contribution is
\[
  \gamma_1(p)=\frac{\lambda}{8\pi^2}E_1(p).
\]

For \(M\) magnons on a long chain, integrability permits an asymptotic
wavefunction built from plane waves in each ordering chamber.  In the
two-magnon chamber \(n_1<n_2\),
\[
  \Psi(n_1,n_2)
  =
  A_{12}\ee^{\ii(p_1n_1+p_2n_2)}
  +
  A_{21}\ee^{\ii(p_2n_1+p_1n_2)} .
\]
Solving the discrete eigenvalue equation at adjacent sites gives
\[
  \frac{A_{21}}{A_{12}}
  =
  S(p_1,p_2),
\]
where, after introducing
\[
  u=\frac12\cot\frac p2,\qquad
  \ee^{\ii p}=\frac{u+\ii/2}{u-\ii/2},
\]
the scattering phase is
\[
  S(u_1,u_2)
  =
  \frac{u_1-u_2+\ii}{u_1-u_2-\ii}.
\]
The periodic closed-chain Bethe equations are
\begin{equation}
  \left(\frac{u_j+\ii/2}{u_j-\ii/2}\right)^L
  =
  \prod_{\substack{k=1\\k\neq j}}^M
  \frac{u_j-u_k+\ii}{u_j-u_k-\ii},
  \qquad
  \prod_{j=1}^M\frac{u_j+\ii/2}{u_j-\ii/2}=1 .
  \label{eq:planar-n4-one-loop-su2-bethe}
\end{equation}
The second equation is the cyclicity condition.  The energy is
\[
  E=\sum_{j=1}^M\frac{1}{u_j^2+1/4}.
\]

The finite-chain transfer-matrix derivation of
\eqref{eq:planar-n4-one-loop-su2-bethe} is exactly the rational \(SU(2)\)
construction in
Chapter~\ref{ch:integrable-algebraic-bethe-ansatz-transfer-matrices}.  The
physical interpretation differs: here the chain is the color trace of a
four-dimensional gauge-theory operator, not the spatial circle of a
two-dimensional relativistic QFT.

\section{Konishi at One Loop}

The shortest non-BPS state in this sector is the length-four two-magnon
descendant represented by \(\operatorname{tr}([X,Z]^2)\).  Set
\[
  L=4,\qquad M=2,\qquad
  u_1=\frac{1}{2\sqrt3},\qquad u_2=-\frac{1}{2\sqrt3}.
\]
Then
\[
  \frac{u_1+\ii/2}{u_1-\ii/2}=\ee^{2\pi\ii/3},
  \qquad
  \frac{u_2+\ii/2}{u_2-\ii/2}=\ee^{-2\pi\ii/3},
\]
so the cyclicity product is \(1\).  The first Bethe equation reads
\[
  \ee^{8\pi\ii/3}
  =
  \frac{u_1-u_2+\ii}{u_1-u_2-\ii}
  =
  \ee^{2\pi\ii/3},
\]
and the second is its inverse.  The spin-chain energy is
\[
  E=\frac{1}{1/12+1/4}+\frac{1}{1/12+1/4}=6 .
\]
Therefore the one-loop anomalous dimension is
\[
  \gamma_{\rm K}^{(1)}
  =
  \frac{\lambda}{8\pi^2}\,6
  =
  \frac{3\lambda}{4\pi^2}.
\]

\section{Beyond One Loop}

At higher loop order the planar dilatation operator becomes a long-range spin
chain.  In the \(SU(2)\) sector its range grows with loop order, and for an
operator of finite length \(L\) interactions can wrap around the trace once
the perturbative range reaches \(L\).  Thus there are two different spectral
problems:
\[
  \text{asymptotic long chain}
  \qquad\text{and}\qquad
  \text{finite trace with wrapping}.
\]
The next chapter describes the asymptotic all-loop Bethe ansatz.  The
following chapter explains how mirror TBA and the Y-system incorporate
finite-length effects.
